\section{Asking for Life, and Given a Word}
\label{sec:opening}

The Twentieth Sunday after Pentecost, \latin{Dominica vigesima post
Pentecosten}, is a Sunday of the second class in the Proper of Time of the 1962
\work{Missale Romanum}. Its Mass begins in the right-hand column of p.~412 of
the typical edition, directly beneath the Postcommunion of the Nineteenth
Sunday, and ends in the right-hand column of p.~413, where the Twenty-first
Sunday begins; its ten elements carry the marginal numbers 1689--1698. In 2026
it falls on 11~October, the second Sunday of the month, a day the calendar gives
to the Maternity of the Blessed Virgin Mary, a feast of the second class. A
Sunday of the second class is preferred to a feast of the same class that is not
a feast of the Lord (\work{Rubricae generales} 16), and the feast is
commemorated wherever the rubrics admit an ordinary commemoration: in every low
Mass and in a conventual Mass (RG 108, 111~b; \work{Rubricae generales Missalis}
434~b). There its Collect, Secret and Postcommunion follow the Sunday's, each
under its own conclusion (RGMR 437~a), and nothing else of its formulary comes
with them, not even its Preface (RGMR 483). A sung Mass that is not conventual,
such as a parish \latin{Missa cantata}, has the Sunday's three orations alone
(RG 111~a; RGMR 434~a). The Mass is said in green (RG 127~b), with Gloria, with
the Credo the rubrics require on every Sunday (RGMR 475~a), and with the Preface
of the Most Holy Trinity, which the Missal prints as a direction after the
Secret and the general rubrics appoint for the Sundays of the second class
outside the Christmas and Easter seasons (RGMR 494~b).

The Mass opens with a confession. Its Introit is cut from the prayer that
Azarias prayed in the furnace at Babylon (Dan 3:24--45): ``Whatever thou hast
done to us, O Lord, thou hast done by a just judgment: for we have sinned and
disobeyed thy commandments: but glorify thy name, and deal with us according to
thy great mercy,'' in the English of the 1861 Philadelphia hand missal. Its
verse is the first of the longest psalm, ``Blessed are the undefiled in the
way, who walk in the law of the Lord.'' The Collect asks pardon and peace of a
God who is appeased, \latin{placatus}, so that the faithful may be cleansed
from every offence and serve him \latin{secura mente}, with a mind at rest. The
Epistle, from the fifth chapter of Ephesians, tells them to walk carefully in
evil days, ``redeeming the time'', to be filled not with wine but with the Holy
Spirit, and to sing and make melody in their hearts, giving thanks always for
all things. The Gradual lifts the eyes of all creatures to the God who gives
them food ``in due season''; the Alleluia answers with a heart that is ready to
sing.

The Gospel is the healing of the ruler's son (Jn 4:46--53). A royal official
whose son lies dying at Capharnaum comes to Jesus and asks him to come down
before the boy dies. Jesus answers, ``Unless you see signs and wonders, you
believe not,'' and then, to the father's second plea, gives him not his
presence but a word: ``Go thy way. Thy son liveth.'' ``The man believed the word
which Jesus said to him and went his way.'' On the road his servants meet him
with the news, he asks the hour, and it was the hour at which the word was
spoken; ``and himself believed, and his whole house.'' The Offertory then sings
the first verse of the exiles' psalm, ``Upon the rivers of Babylon, there we sat
and wept: when we remembered Sion,'' in the Missal's form, which addresses the
city as ``thee''. The Secret asks that the mysteries be heavenly medicine,
\latin{caelestem medicinam}, and purge the vices of the heart. The Communion
reminds God of his word, ``in which thou hast given me hope'', the hope that
``hath comforted me in my humiliation''; and the Postcommunion asks that those
who have received the sacred gifts may always obey the commandments that the
Introit confessed broken.

A father asks for his son's life and is given a word; captives sing by foreign
rivers; a people confesses its sin and asks mercy. Three questions rise from
these texts, and each gives one of the readings that follow. The first asks
what faith Christ asks of the father, and why he heals by a word without going
down. St Gregory the Great answers that the father sought the bodily presence
of one who was absent nowhere, and that the Lord healed by command alone; St
John Chrysostom, that Christ healed the father's mind together with the son's
body, teaching him to attend not to signs but to his word. The second asks how
a people that has sinned and lives in exile is to pray, sing and spend its
time, and it is heard from the Introit, the Offertory and the Epistle. St
Augustine reads the rivers of Babylon as all that is loved here and passes, by
which the citizens of Jerusalem are to sit and weep and not plunge in; St
Hilary of Poitiers reads the psalm as the soul's exile from the heavenly
Jerusalem, its mother, since Adam's sin. The third asks whose story the healing
at Capharnaum tells. St Augustine reads the episode as a figure of the people
of the Jews, who saw signs and scarcely believed, and of the nations, who
believe having seen none; St Thomas Aquinas, as a figure of the Fathers of the
Old Testament praying for a sick people. Every element of the formulary has a
place in each reading, and each closes with four senses of its own, literal,
allegorical, moral and anagogical.

\subsection{The appointed elements at a glance}

\begin{maptable}
Introit & \latin{Omnia, quae fecisti nobis}; Dan 3:31, 29 and 35, as printed;
verse Ps 118:1 (Heb.\ 119) &
Azarias's confession in the furnace: all that God has done was done in true
judgment, ``for we have sinned and disobeyed thy commandments''; then the plea,
glorify thy name and deal with us according to the multitude of thy mercy. The
verse blesses ``the undefiled in the way, who walk in the law of the Lord''.\\
Collect & \latin{Largire, quaesumus, Domine} &
In the 1861 English, ``Favourably grant, we beseech thee, O Lord, thy servants
both pardon and peace; that, being cleansed from the guilt of all their
offences, they may serve thee with secure minds.''\\
Epistle & Eph 5:15--21 &
Walk circumspectly, as wise, ``redeeming the time, because the days are evil'';
understand the will of God; be filled not with wine but with the Holy Spirit;
sing and make melody in your hearts; give thanks always for all things; be
subject one to another in the fear of Christ.\\
Gradual & \latin{Oculi omnium}; Ps 144:15--16 (Heb.\ 145) &
``The eyes of all hope in thee, O Lord: and thou givest them meat in due
season. Thou openest thy hand, and fillest with blessing every living
creature.''\\
Alleluia & \latin{Paratum cor meum}; Ps 107:2 (Heb.\ 108) &
``My heart is ready, O God, my heart is ready''; in the Missal's form the song
is addressed to God, who is named the singer's glory.\\
Gospel & Jn 4:46--53 &
The ruler's son dying at Capharnaum; ``Unless you see signs and wonders, you
believe not''; ``Go thy way. Thy son liveth''; the father believes the word,
goes, learns the hour, and believes with his whole house.\\
Offertory & \latin{Super flumina Babylonis}; Ps 136:1 (Heb.\ 137) &
The captives sit and weep by the rivers of Babylon, remembering Sion, which the
Missal's form addresses as ``thee''.\\
Secret & \latin{Caelestem nobis praebeant} &
In the 1861 English, ``May these mysteries, O Lord, we beseech thee, procure us
a heavenly remedy, and cleanse away the vices of our hearts.''\\
Communion & \latin{Memento verbi tui}; Ps 118:49--50 (Heb.\ 119) &
God is asked to remember his word, ``in which thou hast given me hope. This
hath comforted me in my humiliation''; the psalm's reason, that the word has
given life, is not sung.\\
Postcommunion & \latin{Ut sacris, Domine} &
In the 1861 English, ``That we may be worthy of thy sacred gifts, O Lord:
grant, we beseech thee, we may always obey thy commandments.''\\
\end{maptable}
